\section{Rigidity preservation under Henneberg moves (dimension 2)}
\label{sec:henneberg-rigidity}

This section is the dimension-2 analytical complement of
\cref{sec:henneberg}: both Henneberg moves preserve injective generic
rigidity (\cref{def:isGenericallyRigidInj}) when applied to a
graph that is itself injectively generically rigid in dimension 2.
Combined with the strengthened decomposition
\cref{thm:isLaman-exists-typeI-or-typeII-reverse} and the $K_2$ base
case \cref{thm:top-fin-two-isGenericallyRigidInj}, this is the inductive
step that proves the $\Leftarrow$ direction of Laman's theorem in
\cref{sec:laman-theorem}.

The preservation proofs are rank--nullity arguments on the rigidity
map. For each move, one constructs a kernel-to-kernel linear map via
restriction $x \mapsto x \circ \mathrm{some}$, shows it lands in
$\ker (\mathrm{RigidityMap}\,G\,p)$, and shows it is \emph{injective}
because the new edges through $\mathrm{none}$ pin down the value of
$x$ at $\mathrm{none}$. Injectivity transports the dimension bound
from the underlying-graph kernel to the move-graph kernel. The
specific dimension-2 inputs are:

\begin{itemize}
  \item Type~I: place the fresh vertex at a point $q$ for which
    $q - p_a$ and $q - p_b$ are linearly independent (``off the line
    through $p_a$ and $p_b$''), so that the two new edges pin
    $x(\mathrm{none})$.
  \item Type~II: place $q$ \emph{on} the line through $p_a$ and $p_b$
    (so the two new edges in that direction together recover the
    deleted-edge constraint), with $q - p_a$ and $q - p_c$ linearly
    independent (so the third new edge still pins $x(\mathrm{none})$).
    This forces the input triple $(p_a, p_b, p_c)$ to be
    non-collinear; the unconditional Type~II theorem produces such a
    placement by perturbation, via openness of infinitesimal rigidity
    (\cref{lem:isInfinitesimallyRigid-eventually}).
\end{itemize}

\subsection{Type~I rigidity preservation}

\begin{theorem}[Type~I rank-nullity core]
  \label{thm:typeI-isInfinitesimallyRigid-extend}
  \lean{SimpleGraph.Henneberg.typeI_isInfinitesimallyRigid_extend}
  \leanok
  \uses{def:typeI, def:isInfinitesimallyRigid, lem:rigidityMap-apply}
  Let $V$ be finite, $G$ a graph on $V$, and
  $p : V \to \R^{2}$ a framework that is infinitesimally rigid for
  $G$. Let $a, b \in V$ and $q \in \R^{2}$, and suppose
  $q - p_a$ and $q - p_b$ are linearly independent. Then the extended
  placement $\hat p : \mathrm{Option}\,V \to \R^{2}$ with
  $\hat p(\mathrm{none}) = q$ and $\hat p(\mathrm{some}\,u) = p_u$ is
  infinitesimally rigid for $\mathrm{typeI}\,G\,a\,b$.
\end{theorem}
\begin{proof}
  \leanok
  Define $\rho : \ker (\mathrm{RigidityMap}\,(\mathrm{typeI}\,G\,a\,b)\,\hat p)
  \to_{\R} \ker (\mathrm{RigidityMap}\,G\,p)$ by
  $\rho(x) = x \circ \mathrm{some}$. The map lands in
  $\ker (\mathrm{RigidityMap}\,G\,p)$ because every $G$-edge $\{u, v\}$ lifts to
  $\{\mathrm{some}\,u, \mathrm{some}\,v\}$ in
  $\mathrm{typeI}\,G\,a\,b$ with the same rigidity-row formula.

  $\rho$ is injective: if $\rho(x) = \rho(y)$ then $x$ and $y$ agree
  on $\mathrm{some}\,[V]$. The two new-edge constraints at
  $\{\mathrm{none}, \mathrm{some}\,a\}$ and
  $\{\mathrm{none}, \mathrm{some}\,b\}$ say
  \[
    \langle q - p_a,\, x(\mathrm{none}) - x(\mathrm{some}\,a) \rangle = 0
    \quad\text{and}\quad
    \langle q - p_b,\, x(\mathrm{none}) - x(\mathrm{some}\,b) \rangle = 0
  \]
  (and the same for $y$); subtracting gives that
  $x(\mathrm{none}) - y(\mathrm{none})$ is orthogonal to both
  $q - p_a$ and $q - p_b$. In $\R^{2}$ a vector orthogonal to two
  linearly independent vectors is zero, so $x = y$.

  Injectivity transports the kernel-dimension bound: $\dim
  \ker(\mathrm{RigidityMap}\,(\mathrm{typeI}\,G\,a\,b)\,\hat p) \le
  \dim \ker(\mathrm{RigidityMap}\,G\,p) \le 3$.
\end{proof}

\begin{theorem}[Type~I preserves injective generic rigidity in dimension 2]
  \label{thm:typeI-isGenericallyRigidInj-two}
  \lean{SimpleGraph.Henneberg.typeI_isGenericallyRigidInj_two}
  \leanok
  \uses{def:typeI, def:isGenericallyRigidInj,
    thm:typeI-isInfinitesimallyRigid-extend}
  Let $V$ be finite, $G$ a graph on $V$ that is injectively
  generically rigid in dimension $2$, and $a \ne b$ in $V$. Then
  $\mathrm{typeI}\,G\,a\,b$ is injectively generically rigid in
  dimension $2$.
\end{theorem}
\begin{proof}
  \leanok
  Pick an injective rigid placement $p : V \to \R^{2}$ for $G$.
  Then $p_a \ne p_b$ (by injectivity and $a \ne b$), so the line
  through $p_a$ and $p_b$ is well-defined. The set of $q \in \R^{2}$
  that either lie on this line or coincide with some $p_v$ is a
  union of a one-parameter family and a finite set, hence has empty
  interior; pick $q$ off both. The pair $(q - p_a, q - p_b)$ is
  linearly independent (because $q$ is off the line), so the extended
  placement is infinitesimally rigid
  (\cref{thm:typeI-isInfinitesimallyRigid-extend}). Injectivity of the
  extension follows from injectivity of $p$ and $q \notin p[V]$.
\end{proof}

\subsection{Type~II rigidity preservation}

The Type~II proof has the same shape as the Type~I proof, but places
the new vertex $q$ differently. Place $q$ on the line through $p_a$
and $p_b$, distinct from both.
Then the two new edges $\{\mathrm{none}, \mathrm{some}\,a\}$ and
$\{\mathrm{none}, \mathrm{some}\,b\}$ have rigidity-rows along the
common direction $p_b - p_a$; combining them recovers the
constraint of the \emph{deleted} edge $\{a, b\}$. The third new edge
$\{\mathrm{none}, \mathrm{some}\,c\}$ then pins
$x(\mathrm{none})$, provided $q - p_c$ is independent of the line
direction --- equivalently, $(p_a, p_b, p_c)$ are non-collinear.

\begin{theorem}[Type~II rank-nullity core, non-collinear case]
  \label{thm:typeII-isInfinitesimallyRigid-extend}
  \lean{SimpleGraph.Henneberg.typeII_isInfinitesimallyRigid_extend}
  \leanok
  \uses{def:typeII, def:isInfinitesimallyRigid, lem:rigidityMap-apply}
  Let $V$ be finite, $G$ a graph on $V$, and
  $p : V \to \R^{2}$ infinitesimally rigid for $G$. Let
  $a, b, c \in V$ and $q \in \R^{2}$, and suppose:
  \begin{enumerate}
    \item $q - p_a = \alpha \cdot (p_b - p_a)$ for some
      $\alpha \in \R$ with $\alpha \ne 0, 1$ (i.e.\ $q$ lies on the
      line through $p_a, p_b$, distinct from both); and
    \item $q - p_a$ and $q - p_c$ are linearly independent.
  \end{enumerate}
  Then the extended placement
  $\hat p : \mathrm{Option}\,V \to \R^{2}$
  with $\hat p(\mathrm{none}) = q$ and
  $\hat p(\mathrm{some}\,u) = p_u$ is infinitesimally rigid for
  $\mathrm{typeII}\,G\,a\,b\,c$.
\end{theorem}
\begin{proof}
  \leanok
  As in \cref{thm:typeI-isInfinitesimallyRigid-extend}, define
  $\rho : \ker (\mathrm{RigidityMap}\,(\mathrm{typeII}\,G\,a\,b\,c)\,\hat p)
   \to_{\R} \ker (\mathrm{RigidityMap}\,G\,p)$
  by $\rho(x) = x \circ \mathrm{some}$. We must show that $\rho(x)$ lies
  in $\ker (\mathrm{RigidityMap}\,G\,p)$, that is, its rigidity-map value
  vanishes on every $G$-edge. For $\{u, v\} \ne \{a, b\}$, the edge lifts
  to a $\mathrm{typeII}$-edge with the same row formula, so the
  constraint transports.

  The deleted edge $\{a, b\}$ is recovered from the two new-edge
  constraints at $\{\mathrm{none}, \mathrm{some}\,a\}$ and
  $\{\mathrm{none}, \mathrm{some}\,b\}$. The first hypothesis gives
  $q - p_a = \alpha (p_b - p_a)$ and
  $q - p_b = (\alpha - 1)(p_b - p_a)$ with $\alpha, \alpha - 1 \ne 0$,
  so the two new-edge rows reduce, after dividing by $\alpha$ and
  $\alpha - 1$, to
  \[
    \langle p_b - p_a,\, x(\mathrm{none}) - x(\mathrm{some}\,a) \rangle = 0,
    \quad
    \langle p_b - p_a,\, x(\mathrm{none}) - x(\mathrm{some}\,b) \rangle = 0.
  \]
  Subtracting the second from the first gives
  $\langle p_b - p_a,\, x(\mathrm{some}\,a) - x(\mathrm{some}\,b)
   \rangle = 0$, which is the $\{a, b\}$-row of
   $\mathrm{RigidityMap}\,G\,p$ evaluated at $\rho(x) = x \circ \mathrm{some}$.

  Injectivity of $\rho$ uses the second hypothesis: if $\rho(x) = \rho(y)$
  then $x$ and $y$ agree on $\mathrm{some}\,[V]$, so the new-edge
  constraints at $\{\mathrm{none}, \mathrm{some}\,a\}$ and
  $\{\mathrm{none}, \mathrm{some}\,c\}$ force $x(\mathrm{none}) -
  y(\mathrm{none})$ to be orthogonal to the linearly independent pair
  $(q - p_a, q - p_c)$, hence zero. Rank--nullity then transports the
  kernel-dimension bound as in \cref{thm:typeI-isInfinitesimallyRigid-extend}.
\end{proof}

\begin{theorem}[Type~II preserves injective generic rigidity, non-collinear case]
  \label{thm:typeII-isGenericallyRigidInj-two-of-nonCollinear}
  \lean{SimpleGraph.Henneberg.typeII_isGenericallyRigidInj_two_of_nonCollinear}
  \leanok
  \uses{def:typeII, def:isGenericallyRigidInj,
    thm:typeII-isInfinitesimallyRigid-extend}
  Let $V$ be finite, $G$ a graph on $V$, and
  $p : V \to \R^{2}$ injective and infinitesimally rigid for $G$, with
  $(p_b - p_a, p_c - p_a)$ linearly independent (non-collinear
  triple). Let $a, b, c \in V$ with $a \ne b$. Then
  $\mathrm{typeII}\,G\,a\,b\,c$ is injectively generically rigid in
  dimension $2$.
\end{theorem}
\begin{proof}
  \leanok
  Pick $q$ on the line through $p_a, p_b$, distinct from both and off
  $p[V]$ (all but finitely many points on the line qualify). The
  non-collinearity hypothesis on $(p_a, p_b, p_c)$ ensures
  $q - p_a$ and $q - p_c$ are linearly independent. Apply
  \cref{thm:typeII-isInfinitesimallyRigid-extend}; the
  extended placement is rigid for $\mathrm{typeII}\,G\,a\,b\,c$ and
  injective because $p$ is and $q \notin p[V]$.
\end{proof}

The non-collinearity hypothesis cannot be dropped in
\cref{thm:typeII-isInfinitesimallyRigid-extend}: a collinear placement
of $(p_a, p_b, p_c)$ admits a free one-parameter family of perpendicular
infinitesimal motions of the new vertex, and rigidity fails.
However, openness of infinitesimal rigidity
(\cref{lem:isInfinitesimallyRigid-eventually}) lets us \emph{perturb}
any rigid injective placement to a nearby one with the non-collinearity
property, by moving $p_c$ slightly off the line through $p_a, p_b$.
The result is the unconditional Type~II preservation theorem below.

\begin{theorem}[Type~II preserves injective generic rigidity in dimension 2]
  \label{thm:typeII-isGenericallyRigidInj-two}
  \lean{SimpleGraph.Henneberg.typeII_isGenericallyRigidInj_two}
  \leanok
  \uses{def:typeII, def:isGenericallyRigidInj,
    thm:typeII-isGenericallyRigidInj-two-of-nonCollinear,
    lem:isInfinitesimallyRigid-eventually}
  Let $V$ be finite, $G$ a graph on $V$ that is injectively
  generically rigid in dimension $2$, and $a, b, c \in V$ pairwise
  distinct. Then $\mathrm{typeII}\,G\,a\,b\,c$ is injectively
  generically rigid in dimension $2$.
\end{theorem}
\begin{proof}
  \leanok
  Pick an injective rigid placement $p_0 : V \to \R^{2}$ for $G$.
  If $(p_0(b) - p_0(a), p_0(c) - p_0(a))$ is already linearly
  independent, apply
  \cref{thm:typeII-isGenericallyRigidInj-two-of-nonCollinear} directly.

  Otherwise the triple is collinear:
  $p_0(c) - p_0(a) = \gamma^{-1} \cdot (p_0(b) - p_0(a))$
  for some non-zero $\gamma$. Pick a direction $w \in \R^{2}$ such that
  $w$ and $p_0(b) - p_0(a)$ are linearly independent, and consider the
  one-parameter family of perturbed placements
  $p_t(v) := p_0(v)$ for $v \ne c$ and
  $p_t(c) := p_0(c) + t \cdot w$. The map $t \mapsto p_t$ is
  continuous, $p_0$ is rigid, and rigidity is open in the placement
  (\cref{lem:isInfinitesimallyRigid-eventually}), so the perturbed
  placement remains rigid for $t$ in a neighborhood of $0$.
  Injectivity is also open (each pair $(u, v)$ with $u \ne v$ stays
  separated under continuous perturbation). For any $t \ne 0$ in this
  open neighborhood, the perturbed direction
  $p_t(c) - p_t(a) = \gamma^{-1} (p_0(b) - p_0(a)) + t \cdot w$ is no
  longer in the span of $p_0(b) - p_0(a)$, so the non-collinearity
  condition holds. Apply
  \cref{thm:typeII-isGenericallyRigidInj-two-of-nonCollinear} with
  this perturbed placement.
\end{proof}
